\begin{abstract}
\noindent
Let $\Weyl_n$ be the algebraic Weyl curvature operators on $\R^n$, let $Q(W)=W^2+W^{\#}$ be Hamilton's quadratic map, and let $P(W)=\langle W,Q(W)\rangle$. The problem given here asks that $P$, on the unit sphere of $\Weyl_n$, be maximized on the Weyl tensor of \emph{products of spheres} for $n\ge12$, granting operators diagonal in a suitable basis. It is \textbf{B\"ohm--Wilking Conjecture A}, open. The maximizer is $S^{n/2}\times S^{n/2}$ at $\theta_0(n)=\sqrt{2(n-1)(n-2)}/n$, derived rather than assumed.

\textbf{Problem~P is not solved.} \textbf{The diagonal case is a theorem here, not a hypothesis}: it holds at every $n\ge4$ and both parities, with its equality case (Lemma~57). The published theorem covers even $n\ge16$, and Lemma~1 carries it to $n=12,14$ as a corollary of Lemma~57, which subsumes both. \textbf{At $n=12$ Conjecture~A holds on an exact certificate built outside this project} (Lemma~254), whose conventions are independently checked only in part. The threshold is derived exactly as $n=6+2\sqrt7$ (Lemma~2). A compression of the criticality equation (Lemma~7) splits the critical set by a slack $S\ge0$: the graded half is settled above $\thstar$ for the pair-graded points (Lemma~69), and nondegeneracy across it (Lemmas~27, 28) makes every graded orbit open in the critical set. What remains is one inequality,
\[
\sup\big\{\theta:Q(W)=\theta W,\ S=0,\ W\text{ non-diagonal}\big\}\le\thstar ,
\]
which closes every even $n\ge12$ alone. \textbf{Criticality is part of that statement}: dropped, it is false at every even $n\ge12$, the $\so(6)\oplus\so(6)$-graded subspace of $\Weyl_{12}$ having dimension $834$ against a $54$-dimensional diagonal locus. It is well posed because non-diagonality is $O(n)$-invariant, and \emph{decidable} without a frame search (Lemma~37). \textbf{What is measured is a lower bound}: over the catalogued critical points the supremum equals $\thstar$, attained at $W_{\CP^2}$, so the residual's own direction, $\le$, is untouched by it. \textbf{Diagonality, not $S$, is load-bearing}: $S^5\times S^7$ has slack zero at $1.228239>\thstar$.

Three results place the problem. \textbf{On the compact symmetric spaces the conjecture holds outright} at every $n\ge12$, with its equality case: a dimension count caps every irreducible space at $\thstar$ (Lemma~47) and a partition argument puts the reducible minimum at the balanced product (Lemma~48), the two meeting at $n^2-12n+8\ge0$: the threshold polynomial, recovered. \textbf{$P$ is the left-hand side of B\"ohm--Wilking's cone-invariance criterion}, whose open end demands exactly $\theta_0(n)$ (Lemma~43) and fails for $4\le n\le11$ (Lemma~44): constant and threshold both forced from outside the sphere products. \textbf{And $P$ is a function of two spectra} (Lemma~50): $\theta=\tfrac32p_3(\mu)-\tfrac12p_3(\omega)$, with $\omega$ from the curvature operator of the second kind. Its collapse to one spectrum is Lovelock's identity, holding for $n\le5$ (Lemma~51), which places the only sharp $\thstar$ in the literature. \textbf{No route yet supplies a transfer inequality coupling the two cubics.}

\textbf{Two eigenvalues suffice} (Lemma~105): every critical point with a two-value spectrum has $\theta\le\theta_0(n)$, so Problem~P follows from a two-value maximiser. On the symmetric three-value family that question reduces to \emph{statement H$_\#$}, the sharp norm $(n-2)/\sqrt{\dim\so(n)}$ of Hamilton's $\#$ on symmetric operators (Lemma~259), and H$_\#$ holds on the torus families of $\so(k)$ and $\su(k)$ at every $n\ge14$ (Lemmas~277, 281--283, proved in Lean), on those of $\mathfrak g_2$, $\sype(2)$ and $\sype(3)$, and for $\mathfrak f_4$, $\mathfrak e_6$, $\mathfrak e_7$, $\mathfrak e_8$ in every embedding (Lemmas~286, 287).

Twenty-eight proven negatives (\S\ref{sec:fenceposts}) delimit what cannot work. \textbf{Labels are literal}: a \emph{Lemma} is proved, a \emph{Conjecture} is a candidate, and a result resting on machine computation names its arithmetic in a \emph{Verification} block beside the proof rather than inside it.
\end{abstract}

% ==============================================================================
